\section{Introduction}
\label{sec:intro}

Ask what it costs to run a million agents and the answer usually arrives as one
number. That number is almost always the memory floor of a container, because
the deployment shape that has become idiomatic for agent platforms is one
container per agent, and a container on the commodity serverless platforms
reserves 128\,MB at the low end and 512\,MB at the
high~\cite{modal,e2b,cfworkers}. A million of those is between 128 and 512
terabytes of resident memory, which is a datacenter, and a platform that
believes this about itself will build a scheduler, a billing model and a
market position around a scarcity its own hardware does not impose.

The estimate is wrong because the unit is wrong. ``An agent'' is not one thing
with one cost. It is three things that happen to be discussed together:

\begin{description}[leftmargin=1.4cm,style=nextline]
\item[State] what the agent will resume from --- its identity, its position in
  a workflow, its handle on a conversation. This exists whether or not the
  agent is running, and it is data.
\item[Execution] the context that holds the agent's loop while it waits: a
  stack, a scheduler entry, a place for a continuation to live. This exists
  only while the agent is resident.
\item[Isolation] the boundary the agent's code runs behind, which is what stops
  agent $a$ reading agent $b$'s filesystem or exhausting the host. This exists
  only while the agent's code is actually executing.
\end{description}

Nothing requires these to be the same object, and when they are separated they
turn out to differ by five orders of magnitude. This paper measures each one on
the same laptop, at a fleet of a million, and reports what it found --- including
where the expensive primitive is worth its price.

The finding that motivates the rest: \textbf{an agent that is doing nothing does
not need an execution context or an isolation boundary at all.} It needs a row.
A row costs \rFleetPer{}\,bytes \meas{} and can be brought back into a running
loop in \rFleetResume{}\,ms \meas{}, which is fast enough that keeping the loop
alive to avoid the resume is a false economy for any agent idle longer than a
few milliseconds. The dormant majority of a fleet --- which is most of any real
fleet, most of the time --- therefore costs storage rather than memory, and
storage is four orders of magnitude cheaper.

This is the same argument, one level down, that the companion tenant-residency
paper makes about databases~\cite{unifiedtenant}: there the unit was a whole
per-tenant SQLite file at ${\approx}0.25$\,MB of resident memory, and the result
was that memory scales with concurrency rather than registrations. Here the unit
is a single agent inside such a file, and the result has the same shape: cost
tracks \emph{activity}, not \emph{population}. That the two hold at different
granularities is the point --- it is a property of the decomposition, not of a
particular size.

\subsection*{Contributions}

\begin{enumerate}[leftmargin=*]
\item A separation of the per-agent cost into state, execution and isolation,
  with each measured independently rather than inferred from a deployment shape
  (\S\ref{sec:state}--\S\ref{sec:isolation}).
\item Measured residency for a fleet of \rFleetAgents{} dormant agents:
  \rFleetPer{}\,bytes each on disk, written at \rFleetRate{}/s, resumed in
  \rFleetResume{}\,ms in-process (\S\ref{sec:state}).
\item Measured residency for \rFleetAgents{} \emph{live} execution contexts:
  \rGoPer{}\,bytes of heap per goroutine, reproducible to the byte across five
  runs, with spawn and wake latencies in the hundreds of nanoseconds
  (\S\ref{sec:execution}).
\item Instantiation cost for three in-process language runtimes --- WebAssembly,
  JavaScript and Python --- all in the single-digit to tens of microseconds, and
  all without a container (\S\ref{sec:code}).
\item A like-for-like cold-start comparison of the two isolation primitives that
  are usually conflated, with an argument for why the \rOciRatio{}$\times$ ratio
  between them is not a defect (\S\ref{sec:isolation}).
\item A capacity and cost model built only from the measured constants, and a
  statement of the conditions under which it stops holding
  (\S\ref{sec:capacity}--\S\ref{sec:limits}).
\end{enumerate}
